\subsection{自由 magma 中的 Hall 集}

设 \(\mathbf{X}\) 是一个集合， \(\mathrm{M}(\mathbf{X})\) 是在 \(\mathbf{X}\) 上构造的自由 magma，而 \(\mathrm{M}^n (\mathbf{X})\) （其中 \(n\in \mathbf{N}^*\) ）是 \(\mathrm{M}(\mathbf{X})\) 中长度为 \(n\) 的元素组成的集合 (no. 1)。若 \(w\in \mathrm{M}(\mathbf{X})\) 且 \(l(w)\geqslant 2\) ，用 \(\alpha (w)\) 与 \(\beta (w)\) 表示 \(\mathrm{M}(\mathbf{X})\) 的元素，它们由关系式 \(w = \alpha (w)\beta (w)\) 确定；则 \(l(\alpha (w)) < l(w)\) ， \(l(\beta (w)) < l(w)\) 。最后，对每个 \(u,v\) （属于 \(\mathrm{M}(\mathbf{X})\) ），用 \(u^n v\) 表示对整数 \(m\geqslant 0\) 用归纳法由 \(u^{0}v = v\) 与 \(u^{m + 1}v = u(u^{m}v)\) 定义的元素。

定义 2. 相对于 \(\mathbf{X}\) 的 Hall 集是任何一个带全序的子集 \(\mathbf{H}\) （\(\mathbf{M}(\mathbf{X})\) 中的），满足下列条件：

\begin{enumerate}
\def\labelenumi{(\Alph{enumi})}
\item
  若 \(u \in \mathbf{H}\) 、 \(v \in \mathbf{H}\) 且 \(l(u) < l(v)\) ，则 \(u < v\) 。
\item
  \(\mathbf{X} \subset \mathbf{H}\) ，并且 \(\mathrm{H} \cap \mathrm{M}^2(\mathbf{X})\) 由形如 \(xy\) 的乘积组成，其中 \(x, y\) 属于 \(\mathbf{X}\) 且 \(x < y\) 。
\item
  一个元素 \(w\) （\(\mathbf{M}(\mathbf{X})\) 中的、长度为 \(\geqslant 3\) ）属于 \(\mathbf{H}\) ，当且仅当它具有形式 \(a(bc)\) ，其中 \(a, b, c\) 属于 \(\mathbf{H}, bc \in \mathbf{H}, b \leqslant a < bc\) 且 \(b < c\) 。
\end{enumerate}

命题 11. 存在相对于 X 的 Hall 集。

我们将对整数 \(n\) 用归纳法构造集合 \(\mathbf{H}_n \subset \mathbf{M}^n(\mathbf{X})\) 以及这些集合上的一个全序：

\begin{enumerate}
\def\labelenumi{(\alph{enumi})}
\item
  我们记 \(\mathbf{H}_1 = \mathbf{X}\) ，并赋予它一个全序。
\item
  集合 \(\mathbf{H}_2\) 由形如 \(xy\) 的乘积组成，其中 \(x, y\) 属于 \(\mathbf{X}\) 且 \(x < y\) 。我们赋予它一个全序。
\item
  设 \(n \geqslant 3\) 使得带全序的集合 \(\mathrm{H}_1, \ldots, \mathrm{H}_{n-1}\) 都已定义。集合 \(\mathrm{H}_{n-1}' = \mathrm{H}_1 \cup \dots \cup \mathrm{H}_{n-1}\) 具有一个全序，它在 \(\mathrm{H}_1, \ldots, \mathrm{H}_{n-1}\) 上诱导出给定的序关系，并且 \(w < w'\) （若 \(l(w) < l(w')\) ）。我们定义 \(\mathrm{H}_n\) 为乘积 \(a(bc) \in \mathrm{M}^n(\mathbf{X})\) 组成的集合，其中 \(a, b, c\) 属于 \(\mathrm{H}_{n-1}'\) 且满足关系式 \(bc \in \mathrm{H}_{n-1}', b \leqslant a < bc, b < c\) ，并赋予 \(\mathrm{H}_n\) 一个全序。
\end{enumerate}

我们记 \(\mathbf{H} = \bigcup_{n\geqslant 1}\mathbf{H}_n\) ；我们如下定义地赋予 H 全序： \(w\leqslant w'\) 当且仅当 \(l(w) < l(w')\) ，或者 \(l(w) = l(w') = n\) 且 \(w\leqslant w'\) 属于集合 \(\mathbf{H}_n\) 。立即可见 H 是相对于 X 的 Hall 集。

对 X 的每个子集 S，我们把自由 magma M(S) 与它在 M(X) 中的典范像等同。

命题 12. 设 H 是相对于 X 的一个 Hall 集，并设 x、y 属于 X。

\begin{enumerate}
\def\labelenumi{(\alph{enumi})}
\item
  \(H \cap M(\{x\}) = \{x\}\) 。
\item
  假设 \(x < y\) ，并设 \(d_y\) 是 \(\mathbf{M}(\mathbf{X})\) 到 \(\mathbf{N}\) 中的一个同态，使得 \(d_y(y) = 1\) ，且 \(d_y(z) = 0\) （对 \(z \in \mathbf{X}, z \neq y\) ）。元素 \(w \in \mathbf{H} \cap \mathbf{M}(\{x, y\})\) （满足 \(d_y(w) = 1\) ）的集合由形如 \(x^n y\) 的元素组成，其中 \(n\) 是一个整数 \(\geqslant 0\) 。
\end{enumerate}

由定义 2 (B)， \(x \in \mathrm{H}\) 且 \(\mathrm{H} \cap \mathrm{M}^2(\{x\}) = \varnothing\) 。若 \(w \in \mathrm{H} \cap \mathrm{M}(\{x\})\) ，其中 \(n = l(w) \geqslant 3\) ，则由定义 2 (C)，元素 \(\alpha(w)\) 与 \(\beta(w)\) 也属于 \(\mathrm{H} \cap \mathrm{M}(\{x\})\) 。由此对 \(n\) 立即归纳得出 \(\mathrm{H} \cap \mathrm{M}^n(\{x\}) = \varnothing\) （对 \(n \geqslant 2\) ），从而得到 (a)。

我们现在证明 (b)。由定义 2 (B)， \(y \in \mathsf{H}\) 且 \(xy \in \mathsf{H}\) 。我们对 \(n\) 作归纳证明 \(x^n y \in \mathsf{H}\) （\(n\) 为整数 \(\geqslant 2\) ）。而 \(x^n y = x(x^{n-2}y)\) ，且归纳假设蕴涵 \(x^{n-2}y \in \mathsf{H}\) 。又 \(l(x) < l(x^{n-2}y)\) （对 \(n > 2\) 与 \(x < y\) ），从而在任何情形下都有 \(x < x^{n-2}y\) ；定义 2 的条件 (C) 表明 \(x^n y \in \mathsf{H}\) 。另一方面，显然 \(d_y(x^n y) = 1\) 。反之，设 \(w \in \mathsf{H} \cap \mathrm{M}(\{x, y\})\) ，其中 \(d_y(w) = 1\) 。若 \(l(w) = 1\) ，则 \(w = y\) ；若 \(l(w) = 2\) ，则由定义 2 (B) 得 \(w = xy\) 。若 \(l(w) \geqslant 3\) ，则 \(w = a(bc)\) ，其中 \(a, b, c, bc\) 属于 \(\mathsf{H} \cap \mathrm{M}(\{x, y\})\) (定义 2 (C))。 \(d_y(bc) = 0\) 是不可能的，因为这将蕴涵 \(bc \in \mathrm{M}(\{x\})\) ，而由 (a) 这是不可能的。因此 \(d_y(bc) = 1\) 且 \(d_y(a) = 0\) ，从而由 (a) 得 \(a = x\) 。于是对 \(n = l(w)\) 立即归纳得出 \(w = x^{n-2}y\) ，这就完成了 (b) 的证明。

推论. 若 \(\operatorname{Card} X \geqslant 2\) ，则 \(H \cap M^n(X) \neq \varnothing\) （对每个整数 \(n \geqslant 1\) ）。

命题 13. 设 X 是至少含有两个元素的有限集。设 H 表示相对于 X 的一个 Hall 集。那么存在一个 N 到 H 上的严格递增双射 \(p \mapsto w_p\) 以及 H 的子集的一个序列 \((\mathbf{P}_p)_{p \in \mathbf{N}}\) ，它们具有下列性质：

\begin{enumerate}
\def\labelenumi{(\alph{enumi})}
\item
  \(P_{0} = X.\)
\item
  对每个整数 \(p \geqslant 0\) ， \(w_{p} \in \mathbf{P}_{p}\) 。
\item
  对每个整数 \(n \geqslant 1\) ，存在一个整数 \(p(n)\) ，使得 \(\mathbf{P}_p\) 的每个元素都具有长度 \(> n\) （对一切 \(p \geqslant p(n)\) ）。
\item
  对每个整数 \(p \geqslant 0\) ，集合 \(\mathrm{P}_{p+1}\) 由形如 \(w_p^i w\) 的元素组成，其中 \(i \geqslant 0, w \in \mathrm{P}_p\) 且 \(w \neq w_p\) 。
\end{enumerate}

由于 X 是有限的，每个集合 \(\mathrm{M}^{n}(\mathrm{X})\) 都是有限的。令 \(H_{n}=H\cap M^{n}(\mathbf{X})\) （对一切 \(n\geqslant1\) ）。命题 12 的推论表明有限集 \(H_{n}\) 非空。设 \(u_{n}\) 是 \(H_{n}\) 的基数；令 \(v_{0}=0\) ，并令 \(v_{n}=u_{1}+\cdots+u_{n}\) （对 \(n\geqslant1\) ）。由于 \(H_{n}\) 是一个带全序的有限集，存在一个严格递增双射 \(p\mapsto w_{p}\) ，从 N 的区间 \([v_{n-1},v_{n}-1]\) 到 \(H_{n}\) 上。立即可见 \(p\mapsto w_{p}\) 是 N 到 H 上的一个严格递增双射。

设 \(\mathbf{P}_0 = \mathbf{X}\) ，并且对每个整数 \(p \geqslant 1\) ，令 \(\mathbf{P}_p\) 为元素 \(w\) 组成的集合，这些元素属于 \(\mathbf{H}\) ，满足 \(w \geqslant w_p\) ，并且或者 \(w \in \mathbf{X}\) 、或者 \(\alpha(w) < w_p\) （注意，若 \(w\) 的长度为 \(\geqslant2\) ，则由定义 2 的条件 (C)，关系式 \(w\in H\) 蕴涵 \(\alpha(w)\in H\) ）。那么 \(w_{p}\in P_{p}\) ；当 \(w_{p}\in X\) 时这是显然的，而由不等式 \(l(\alpha(w_{p}))<l(w_{p})\) 与定义 2 的条件 (A) 可得（当 \(w_{p}\notin X\) 时）。

因此条件 (a) 与 (b) 都满足。

设 \(n\) 是一个整数 \(\geqslant 1\) ，并设 \(p \geqslant v_n\) 。对一切 \(w \in \mathrm{P}_p\) ，\(l(w) \geqslant l(w_p) > n\) 由映射 \(p \mapsto w_p\) 的定义本身即得。这就证明了 (c)。

我们现在证明，每个形如 \(u = w_{p}^{i}w\) 且满足 \(i \geqslant 0\) 、 \(w \in \mathrm{P}_{p}\) 与 \(w \neq w_{p}\) 的元素都属于 \(\mathrm{P}_{p+1}\) 。若 \(i \neq 0\) ，则 \(l(u) > l(w_{p})\) ，从而 \(u > w_{p}\) 且 \(u \geqslant w_{p+1}\) ；于是 \(u \notin X\) 且 \(\alpha(u) = w_{p} < w_{p+1}\) ，从而 \(u \in \mathrm{P}_{p+1}\) 。若 \(i = 0\) ，则 \(u \in \mathrm{P}_{p}\) 且 \(u \neq w_{p}\) ；于是 \(u > w_{p}\) ，从而 \(u \geqslant w_{p+1}\) ；若 \(u\) 不属于 \(X\) ，则 \(\alpha(w) < w_{p}\) ，从而 \(\alpha(w) < w_{p+1}\) ；于是再次有 \(u \in \mathrm{P}_{p+1}\) 。反之，设 \(u \in \mathrm{P}_{p+1}\) 。我们区分两种情形：

\((\alpha)\) 不存在元素 \(v\) （属于 \(\mathbf{M}(\mathbf{X})\) ）满足 \(u = w_{p}v\) 。由 \(\mathrm{P}_{p + 1}\) 的定义， \(u > w_{p}\) 。此外，若 \(u\notin \mathbf{X}\) ，则由所给假设有 \(\alpha (u)\neq w_p\) ，并且 \(\alpha (u) < w_{p + 1}\) （由于 \(u\in \mathrm{P}_{p + 1}\) ）；因此 \(\alpha (u) < w_{p}\) 。于是 \(u\in \mathrm{P}_p\) 且 \(u\neq w_{p}\) 。

(β) 存在 \(v\) 属于 \(\mathbf{M}(\mathbf{X})\) ，使得 \(u = w_{p}v\) 。由定义 2，必然或者 \(w_{p} \in \mathbf{X}\) 、 \(v \in \mathbf{X}\) 且 \(w_{p} < v\) ，或者 \(v \notin \mathbf{X}\) 且 \(\alpha(v) \leqslant w_{p} < v\) 。在两种情形下都有 \(v \in \mathbf{P}_{p+1}\) 。

于是存在一个整数 \(i \geqslant 0\) 与一个元素 \(w\) （属于 \(\mathrm{M}(\mathrm{X})\) ），使得 \(u = w_p^i w\) ，并且或者 \(w \in \mathrm{X}\) ，或者 \(w \notin \mathrm{X}\) 且 \(\alpha(w) \neq w_p\) 。若 \(i = 0\) ，则得到上面的情形 \((\alpha)\) ，从而 \(w \in \mathrm{P}_p\) 且 \(w \neq w_p\) 。若 \(i > 0\) ，则上面 \((\beta)\) 的证明对 \(i\) 作归纳便建立起关系式 \(w \in \mathrm{P}_{p+1}\) 与 \(w \neq w_p\) 。假设 \(w \notin \mathrm{X}\) ；由 \(w \in \mathrm{P}_{p+1}\) 推出 \(\alpha(w) \leqslant w_p\) ，并且由于 \(\alpha(w) \neq w_p\) ，我们断定 \(w \in \mathrm{P}_p\) 。这就完成了 (d) 的证明。

例. 假设 X 有两个元素 x、y；设 X 的序使 x < y。命题 11 的证明中所给出的构造给出一个集合 H，它含有 14 个长度为 \(\leqslant 5\) 的元素，列于下表：

\[
\begin{array}{l l l l} \mathrm{H} _ {1} & w _ {1} = x & w _ {2} = y \\ \mathrm{H} _ {2} & w _ {3} = (x y) \\ \mathrm{H} _ {3} & w _ {4} = (x (x y)) & w _ {5} = (y (x y)) \\ \mathrm{H} _ {4} & w _ {6} = (x (x (x y))) & w _ {7} = (y (x (x y))) & w _ {8} = (y (y (x y))) \\ \mathrm{H} _ {5} & w _ {9} = (x (x (x (x y)))) & w _ {10} = (y (x (x (x y)))) & w _ {11} = (y (y (x (x y)))) \\ & w _ {12} = (y (y (x y))) & w _ {13} = ((x y) (x (x y))) & w _ {14} = ((x y) (y (x y))), \end{array}
\]

(H 的元素已按照在每个 \(H_{n}\) 上选取的全序编号)。

